% ECONOMICS_PROBLEM: {"id": "C90-604", "title": "Laboratory-to-field transfer", "jel_code": "C90", "source_id": "OP-0306"}
% !TeX program = xelatex
\documentclass[11pt,a4paper]{article}
\usepackage[margin=23mm]{geometry}
\usepackage{fontspec}
\setmainfont{Latin Modern Roman}
\usepackage{amsmath,amssymb,mathtools}
\usepackage{xcolor,enumitem,booktabs,longtable,array,xurl,hyperref}
\definecolor{muted}{HTML}{53636C}
\hypersetup{colorlinks=true,urlcolor=blue,linkcolor=blue}
\setlength{\parindent}{0pt}
\setlength{\parskip}{5pt}
\newcommand{\R}{\mathbb R}
\newcommand{\E}{\mathbb E}
\newcommand{\N}{\mathbb N}
\newcommand{\ind}{\mathbf 1}
\newcommand{\Var}{\operatorname{Var}}
\newcommand{\Cov}{\operatorname{Cov}}
\newcommand{\supp}{\operatorname{supp}}
\DeclareMathOperator*{\argmin}{arg\,min}
\DeclareMathOperator*{\argmax}{arg\,max}
\newcommand{\entrymeta}[3]{{\sffamily\small\color{muted}\textbf{\texttt{#1}}\quad|\quad #2\par #3\par}}
\newcommand{\Problemref}[1]{\texttt{#1}}
\newcommand{\JELref}[2]{\texttt{#2} (JEL \texttt{#1})}
\begin{document}
\section*{Laboratory-to-field transfer}
\textbf{Problem ID:} \texttt{C90-604}\quad \textbf{JEL:} \texttt{C90}\par
% BEGIN ECONOMICS STATEMENT
\label{problem:OP-0306}
{\sffamily\small\color{muted}Original subject group: Behavioral and experimental economics\par}
\label{definitions:problem:30007588}
\textbf{Audit classification:} Published research agenda. Levitt–List 2007 explicitly raises generalizability and calls for models connecting settings. The potential-outcome transport formula and sensitivity classes are editorial.
\subsection*{Definitions and precise research target}
\textbf{Model and research target.} Let \(S\in\{L,F\}\) indicate a laboratory sample or a target field population. For a specified intervention \(D\in\{0,1\}\), let \(Y_i(d,S)\) denote individual \(i\)'s potential outcome under treatment \(d\) in setting \(S\). Let \(X_i\) collect participant characteristics and \(C_S\) collect contextual features, including scrutiny, stakes, framing, and institutional rules. The field target and a candidate laboratory transport formula are
\[\tau_F=\mathbb E_F[Y_i(1,F)-Y_i(0,F)],\qquad \widetilde\tau_F=\int\mathbb E[Y_i(1,L)-Y_i(0,L)\mid X_i=x]\,dP_F(x).\]
The formula requires overlap in \(X\), internal identification of the laboratory contrast, and invariance of conditional treatment effects across contexts. It therefore does not follow merely from random assignment inside the laboratory. Determine which observed contextual moderators permit a credible extended transport model \(\tau(x,C)\), and quantify errors when context or selection variables are unmeasured. Use deliberately varied laboratory settings and independent field comparisons to evaluate transported predictions, with the treatment and outcome kept conceptually comparable. Distinguish differences in effect magnitudes from different outcome measurement or different versions of the intervention.

\emph{Definitions and maintained assumptions (editorial unless a source definition is named).} Use the same conceptual intervention and outcome units in both settings. \(Y(d,s)\) is a potential outcome under treatment version \(d\) and setting \(s\); consistency gives observed \(Y=Y(D,S)\). Lab assignment satisfies \(D\perp(Y(0,L),Y(1,L))\mid X,S=L\), with \(0<\Pr(D=1\mid X,S=L)<1\). Transport additionally requires \(P_F\ll P_L\) on covariates and equality of conditional mean effects \(\tau_F(x)=\tau_L(x)\) almost everywhere under \(P_F\). Then the displayed formula identifies \(\tau_F\) from lab outcomes and field covariates. This identification implication is established by these assumptions and is not itself an unresolved theorem. If \(|\tau_F(x)-\tau_L(x)|\le\Gamma(x)\), the transport bias is bounded by \(\int\Gamma(x)dP_F(x)\); the research problem is credible calibration of \(\Gamma\) and moderator selection. A context variable constant within one setting cannot be disentangled from setting without additional varied environments.
\subsection*{Variants and assumption changes}
\begin{enumerate}
\item \textbf{Conditional transport (editorial).} Collect multiple contexts with overlapping scrutiny, stakes and participant covariates. Specify an invariance model \(\tau(x,C)\) and validate its prediction in an independent field setting.
\item \textbf{Sensitivity (editorial).} Bound unmeasured context differences by a justified \(\Gamma(x)\) and incomplete covariate overlap by bounded outcomes. Report the resulting sharp or conservative effect interval instead of extrapolating a point estimate.
\end{enumerate}
\subsection*{References and source audit}
\begin{enumerate}
\item Introduction pp.153–154; discussion pp.170–171. Supports the stated research area; exact displayed specification is editorial. \url{https://bpb-us-w2.wpmucdn.com/voices.uchicago.edu/dist/f/1276/files/2018/09/27-101-25264884-1gx76fi.pdf}
\end{enumerate}
\begin{enumerate}\raggedright
\item\hypertarget{definitions:source:30007588:1}{} Steven D. Levitt; John A. List. 2007. \emph{What Do Laboratory Experiments Measuring Social Preferences Reveal About the Real World?}. Introduction, pp. 153–154; Stakes, p. 165; Conclusions, pp. 170–171 (PDF pp. 1–2, 13, 18–19)..
\url{https://bpb-us-w2.wpmucdn.com/voices.uchicago.edu/dist/f/1276/files/2018/09/27-101-25264884-1gx76fi.pdf}.
DOI: \url{https://doi.org/10.1257/jep.21.2.153}.
\end{enumerate}

\subsection*{Common conventions: Source mathematical conventions}

These conventions apply to each entry unless explicitly replaced.

\textbf{Models and identification.} A primitive model \(m\) specifies all probability laws, feasible choices, information, equilibrium and measurement rules used by its target. The admissible class \(\mathcal M\) is fixed independently of outcomes. If \(P_O(m)\) is its observable probability law and \(\tau(m)\) the target, its identified set at observable law \(P\) is
\[
 \mathcal I_\tau(P)=\{\tau(m):m\in\mathcal M,\ P_O(m)=P\}.
\]
Point identification means that this set is a singleton. An empty set signals incompatibility of data and assumptions. Coordinate bounds need not describe the joint set. Bounds are sharp only if every retained target is attainable or arbitrarily approachable by compatible models. Finite bounds require explicit restrictions on unobserved quantities; unrestricted confounding, hidden wealth, extrapolation or arbitrary equilibrium selection can yield uninformative sets.

\textbf{Causal interventions.} A potential outcome \(Y_i(p)\) is the outcome under a complete policy regime \(p\), including timing, financing, information and market spillovers. The target population and units are fixed before treatment. With interference, use the complete assignment vector or a justified exposure mapping; individual no-interference notation alone cannot identify an economy-wide intervention. A difference of mean potential outcomes does not require the cross-world joint distribution. Conditioning on post-treatment migration, survival, enrollment or employment changes the estimand and needs separate justification.

\textbf{Statistical statements.} An estimator, confidence set, forecast or test specifies a sampling law, model class and loss/coverage criterion. Uniform \((1-\alpha)\)-coverage means \(\inf_{m\in\mathcal M}\Pr_m\{\tau(m)\in C_n\}\geq1-\alpha\), or an explicitly asymptotic version. The whole parameter space trivially has coverage, so informativeness requires a stated width, risk or power objective. A forecasting improvement is relative to a named benchmark, information set and evaluation population.

\textbf{Equilibrium and optimization.} An equilibrium comprises feasible plans, optimality, beliefs where needed, budget constraints and all market/resource clearing. The primitives or a stated benchmark define each correspondence \(\mathcal E(p;m)\). Nonemptiness is assumed or proved, never inferred just from compact policy sets. A selected-equilibrium target uses a declared selection rule; a robust target quantifies over all permitted equilibria. Use suprema or infima unless attainment is established. An \(\varepsilon\)-optimal policy is feasible and within \(\varepsilon\) of the finite optimum value. Report an empty feasible set separately.

\textbf{Welfare and accounting.} Normative utility, interpersonal normalization, social weights, numeraire, discounting and the use of public receipts are primitives. Behavioral utility need not coincide with the welfare criterion. Household welfare includes ownership income and rents once; adding profits again to compensating variation generally double-counts them. A monetary equivalent is defined by an explicit transfer and price convention, with monotonicity and range conditions. Average effects, mechanism contributions, transfers and welfare losses are different targets. Infinite sums require convergence and dynamic feasible plans require terminal or no-Ponzi conditions.

\textbf{Source versions.} A working-paper date, revision date and journal publication year are distinguished. A DOI for a journal version is not automatically the DOI of an earlier manuscript. Locators refer to the stated version and printed pages where specified. A metadata-only check does not verify a claim about a full-text conclusion. Every entry's source subsection narrows the provenance claim accordingly.
% END ECONOMICS STATEMENT
\end{document}
